\chapter{1959:Jaroslav Heyrovsky}

\label{chap:chemistry1959}

The Nobel Prize in Chemistry for the year 1959 was awarded to Jaroslav Heyrovsky.

\textbf{Motivation:}

for his discovery and development of the polarographic methods of analysis

\section{Jaroslav Heyrovsky}

\subsection*{Early Life and Education}

Jaroslav Heyrovsky was born on 1890-12-20 in Prague.

\subsection*{Career and Major Research}

Heyrovsky studied at Charles University in Prague and at University College London, where he earned his doctorate in 1918 under the physical chemist Frederick G. Donnan. He returned to Prague and spent his career at Charles University, becoming professor of physical chemistry. In 1922 he invented polarography, an analytical method in which current is measured as a function of the voltage applied to a dropping mercury electrode. He developed the method into a general technique of electroanalysis, and in 1950 he founded the Polarographic Institute of the Czechoslovak Academy of Sciences in Prague.

\subsection*{Nobel Prize-winning Work}

The work for which Jaroslav Heyrovsky was awarded the Nobel Prize in Chemistry in 1959 was described by the Nobel Committee as: "for his discovery and development of the polarographic methods of analysis".

Polarography records the current-voltage curves obtained when solutions are electrolyzed with a dropping mercury electrode, and the position and height of the resulting waves identify the substance and measure its concentration. The method made possible the quantitative analysis of reducible and oxidizable substances down to very small concentrations.

\subsection*{Significance and Impact}

Polarography became one of the most widely used techniques of analytical chemistry, applied to trace metals, organic compounds, and biological materials in industry and medicine. It stimulated the development of the whole field of modern electroanalytical chemistry.

\subsection*{Later Career and Honors}

Heyrovsky remained in Prague and directed the Polarographic Institute until his death. He died on 1967-03-27 in Prague.

\subsection*{Legacy}

Heyrovsky is regarded as the founder of polarography and of modern electroanalysis, and the Jaroslav Heyrovsky Institute of Physical Chemistry in Prague continues the school he created.

\subsection*{Selected Publications}

\begin{itemize}

\item \emph{A Polarographic Study of the Electro-Kinetic Phenomena of Absorption, Electro-Reduction and Overpotential Displayed at the Dropping Mercury Cathode}, 1934.

\end{itemize}

\clearpage

\section*{Chapter Summary}

The 1959 prize recognized the discovery and development of polarography. Jaroslav Heyrovsky invented the method at Charles University in Prague in 1922, using the dropping mercury electrode to obtain current-voltage curves, and developed it into a widely used technique of analytical chemistry.

